\section{Miary zewnętrzne}
\label{sec:miary-zewnetrzne}

\begin{uwaga}[Rodzaje przedziałów w pokryciach]
\label{uw:rodzaje-pokryc}
W definicji~\ref{def:miara-zewnetrzna-lebesgue} otrzymamy tę samą
wartość $\lambda^{\ast}(A)$, jeżeli dopuścimy wyłącznie przedziały
otwarte albo wyłącznie przedziały domknięte. Możemy też dopuścić
zbiór pusty, przyjmując $|\varnothing|=0$.
\end{uwaga}

\begin{proof}
Ograniczenie rodziny dopuszczalnych pokryć może tylko zwiększyć infimum.
Jeżeli $\lambda^{\ast}(A)=+\infty$, wszystkie rozważane infima są więc
równe $+\infty$. W przeciwnym przypadku ustalmy $\varepsilon>0$
i wybierzmy pokrycie $A$ przedziałami $I_n$ takie, że
\[
  \sum_{n=1}^{\infty}|I_n|<\lambda^{\ast}(A)+\varepsilon.
\]
Każdy $I_n$ ma skończoną długość. Możemy go zawrzeć w przedziale
otwartym $J_n$ o długości mniejszej niż $|I_n|+\varepsilon/2^n$.
Wtedy $J_n$ pokrywają $A$ i
\[
  \sum_{n=1}^{\infty}|J_n|<\lambda^{\ast}(A)+2\varepsilon.
\]
Infimum dla pokryć otwartych jest zatem równe $\lambda^{\ast}(A)$.
Zastąpienie każdego $I_n$ jego domknięciem nie zmienia długości,
więc równość zachodzi również dla pokryć domkniętych.
Puste składniki pokrycia można zastąpić przedziałami jednopunktowymi
długości zero; ich dopuszczenie także nie zmienia infimum.
\end{proof}

\begin{samepage}
\begin{definicja}[Miara zewnętrzna]
\label{def:miara-zewnetrzna}
Miarą zewnętrzną na zbiorze $X$ nazywamy funkcję zbioru
$\mu^{\ast}\colon\mathcal{P}(X)\to[0,+\infty]$, która spełnia warunki:
\begin{enumerate}
  \item $\mu^{\ast}(\varnothing)=0$;
  \item jeżeli $A\subseteq B\subseteq X$, to
    $\mu^{\ast}(A)\leq\mu^{\ast}(B)$;
  \item dla dowolnych $A_n\subseteq X$ zachodzi
    \[
      \mu^{\ast}\left(\bigcup_{n=1}^{\infty}A_n\right)
      \leq\sum_{n=1}^{\infty}\mu^{\ast}(A_n).
    \]
\end{enumerate}
Są to kolejno: zerowanie na zbiorze pustym, monotoniczność
i przeliczalna podaddytywność.
\end{definicja}
\end{samepage}

\begin{uwaga}
Funkcja $\lambda^{\ast}$ jest miarą zewnętrzną na $\R$, zgodnie
z faktami~\ref{fakt:zewnetrzna-zero}--\ref{fakt:zewnetrzna-podaddytywnosc}.
Miara zewnętrzna może przyjmować wartość $+\infty$.
\end{uwaga}

\section{Warunek mierzalności Carathéodory'ego}
\label{sec:caratheodory}

\begin{definicja}[Zbiór mierzalny względem miary zewnętrznej]
\label{def:caratheodory}
Niech $\mu^{\ast}$ będzie miarą zewnętrzną na $X$.
Zbiór $E\subseteq X$ nazywamy $\mu^{\ast}$-mierzalnym, jeżeli
dla każdego zbioru $A\subseteq X$ zachodzi
\begin{equation}
  \mu^{\ast}(A)
  =\mu^{\ast}(A\cap E)+\mu^{\ast}(A\cap E^c).
  \label{eq:caratheodory}
\end{equation}
Dopełnienie $E^c$ jest brane względem $X$.
Równość~\eqref{eq:caratheodory} nazywamy warunkiem mierzalności
Carathéodory'ego.
\end{definicja}

\begin{twierdzenie}[Carathéodory'ego]
\label{tw:caratheodory}
Rodzina $\mathcal{F}$ wszystkich zbiorów $\mu^{\ast}$-mierzalnych
jest sigma-ciałem. Obcięcie
\[
  \mu=\mu^{\ast}|_{\mathcal{F}}
  \colon\mathcal{F}\longrightarrow[0,+\infty]
\]
jest miarą zupełną.
\end{twierdzenie}

\begin{proof}
\emph{Sigma-ciało.}
Zbiór pusty jest mierzalny, bo
$\mu^{\ast}(A\cap\varnothing)+\mu^{\ast}(A\cap X)=\mu^{\ast}(A)$.
Warunek~\eqref{eq:caratheodory} jest symetryczny względem $E$ i $E^c$,
więc $\mathcal{F}$ jest zamknięta na dopełnienia.

Niech $E,F\in\mathcal{F}$. Dzieląc najpierw $A$ przez $E$, a potem
$A\cap E^c$ przez $F$, otrzymujemy
\begin{align*}
  \mu^{\ast}(A)
    &=\mu^{\ast}(A\cap E)+\mu^{\ast}(A\cap E^c\cap F)
      +\mu^{\ast}(A\cap E^c\cap F^c)\\
    &\geq\mu^{\ast}\bigl(A\cap(E\cup F)\bigr)
      +\mu^{\ast}\bigl(A\cap(E\cup F)^c\bigr).
\end{align*}
W nierówności zastosowaliśmy podaddytywność do dwóch pierwszych
składników. Przeciwna nierówność również wynika z podaddytywności,
zatem $E\cup F\in\mathcal{F}$.

Niech teraz $E_n\in\mathcal{F}$ będą parami rozłączne, a
$S_N=\bigcup_{n=1}^{N}E_n$ i $S=\bigcup_{n=1}^{\infty}E_n$.
Kolejne zastosowania warunku mierzalności dają
\[
  \mu^{\ast}(A)
  =\sum_{n=1}^{N}\mu^{\ast}(A\cap E_n)+\mu^{\ast}(A\setminus S_N)
  \geq\sum_{n=1}^{N}\mu^{\ast}(A\cap E_n)
      +\mu^{\ast}(A\setminus S).
\]
Przechodząc z $N$ do nieskończoności i stosując podaddytywność,
otrzymujemy
\[
  \mu^{\ast}(A)
  \geq\sum_{n=1}^{\infty}\mu^{\ast}(A\cap E_n)
       +\mu^{\ast}(A\setminus S)
  \geq\mu^{\ast}(A\cap S)+\mu^{\ast}(A\cap S^c).
\]
Przeciwna nierówność jest automatyczna, więc $S\in\mathcal{F}$.
Dla dowolnego ciągu $F_n\in\mathcal{F}$ stosujemy rozłącznienie
\[
  E_1=F_1,\qquad
  E_n=F_n\setminus\bigcup_{k=1}^{n-1}F_k\quad(n\geq2).
\]
Zamknięcie na dopełnienia i skończone sumy zapewnia $E_n\in\mathcal{F}$,
a równość sum obu ciągów daje $\bigcup_nF_n\in\mathcal{F}$.
Zatem $\mathcal{F}$ jest sigma-ciałem.

\emph{Miara.}
Mamy $\mu(\varnothing)=0$.
Dla rozłącznych $E,F\in\mathcal{F}$, stosując mierzalność $F$ do
$A=E\cup F$, otrzymujemy
\[
  \mu(E\cup F)
  =\mu^{\ast}\bigl((E\cup F)\cap F\bigr)
   +\mu^{\ast}\bigl((E\cup F)\cap F^c\bigr)
  =\mu(F)+\mu(E).
\]
Indukcyjnie dostajemy addytywność dla dowolnej skończonej rodziny.
Dla parami rozłącznych $E_n\in\mathcal{F}$ monotoniczność daje
\[
  \mu\left(\bigcup_{n=1}^{\infty}E_n\right)
  \geq\mu\left(\bigcup_{n=1}^{N}E_n\right)
  =\sum_{n=1}^{N}\mu(E_n).
\]
Po przejściu do granicy otrzymujemy nierówność z sumą nieskończoną.
Nierówność przeciwna wynika z podaddytywności $\mu^{\ast}$,
co dowodzi przeliczalnej addytywności $\mu$.

\emph{Zupełność.}
Niech $Z\subseteq B$, gdzie $B\in\mathcal{F}$ i $\mu(B)=0$.
Z monotoniczności $\mu^{\ast}(Z)=0$, więc dla każdego $A\subseteq X$
mamy $\mu^{\ast}(A\cap Z)=0$ oraz
\[
  \mu^{\ast}(A\cap Z)+\mu^{\ast}(A\cap Z^c)
  =\mu^{\ast}(A\cap Z^c)\leq\mu^{\ast}(A).
\]
Przeciwna nierówność wynika z podaddytywności.
Zatem $Z$ jest mierzalny, czyli $Z\in\mathcal{F}$.
\end{proof}

\begin{wniosek}[Miara Lebesgue'a]
\label{wn:miara-lebesgue}
Niech $\mathcal{L}$ oznacza rodzinę wszystkich
$\lambda^{\ast}$-mierzalnych podzbiorów $\R$. Wówczas
\[
  \lambda=\lambda^{\ast}|_{\mathcal{L}}
\]
jest miarą zupełną, a $(\R,\mathcal{L},\lambda)$ jest zupełną
przestrzenią z miarą. Miarę $\lambda$ nazywamy miarą Lebesgue'a,
a $\mathcal{L}$ --- sigma-ciałem zbiorów mierzalnych w sensie Lebesgue'a.
\end{wniosek}

\begin{uwaga}
Zbiór Vitaliego $V$ z przykładu~\ref{prz:vitali} nie należy do
$\mathcal{L}$. Gdyby był mierzalny, mierzalne byłyby też jego
przesunięcia (jak wykażemy poniżej), a addytywność $\lambda$
prowadziłaby do otrzymanej tam sprzeczności.
\end{uwaga}

\section{Własności miary Lebesgue'a}
\label{sec:wlasnosci-lebesgue}

\begin{samepage}
\begin{fakt}
\label{fakt:wlasnosci-lebesgue}
Miara Lebesgue'a i jej miara zewnętrzna mają następujące własności:
\begin{enumerate}
  \item Jeżeli $E\in\mathcal{L}$ i $x\in\R$, to $E+x\in\mathcal{L}$
    i $\lambda(E+x)=\lambda(E)$.
  \item Każdy zbiór przeliczalny $F\subseteq\R$ należy do
    $\mathcal{L}$ i ma miarę zero. W szczególności
    $\lambda(\{x\})=0$ dla $x\in\R$.
  \item Dla każdego przedziału $I\subseteq\R$ zachodzi
    $\lambda^{\ast}(I)=|I|$.
  \item Każda półprosta $(-\infty,b)$, gdzie $b\in\R$,
    należy do $\mathcal{L}$.
  \item $\mathcal{B}(\R)\subseteq\mathcal{L}$.
\end{enumerate}
\end{fakt}
\end{samepage}

\begin{proof}
\emph{Ad 1.}
Niech $E\in\mathcal{L}$. Dla dowolnego $A\subseteq\R$ stosujemy
mierzalność $E$ do $A-x$. Z niezmienniczości $\lambda^{\ast}$
na przesunięcia wynika
\begin{align*}
  \lambda^{\ast}(A)
    &=\lambda^{\ast}(A-x)\\
    &=\lambda^{\ast}\bigl((A-x)\cap E\bigr)
      +\lambda^{\ast}\bigl((A-x)\cap E^c\bigr)\\
    &=\lambda^{\ast}\bigl(A\cap(E+x)\bigr)
      +\lambda^{\ast}\bigl(A\cap(E+x)^c\bigr).
\end{align*}
Zatem $E+x\in\mathcal{L}$ i
$\lambda(E+x)=\lambda^{\ast}(E+x)=\lambda^{\ast}(E)=\lambda(E)$.

\emph{Ad 2.}
Przedział $[x,x]$ pokrywa $\{x\}$ i ma długość zero,
więc $\lambda^{\ast}(\{x\})=0$.
Przeliczalna podaddytywność daje $\lambda^{\ast}(F)=0$
dla każdego przeliczalnego $F$.
Każdy zbiór miary zewnętrznej zero jest mierzalny:
dla dowolnego $A$ jeden składnik warunku Carathéodory'ego jest zerowy,
a drugi nie przekracza $\lambda^{\ast}(A)$; przeciwną nierówność
daje podaddytywność. Stąd $F\in\mathcal{L}$ i $\lambda(F)=0$.

\emph{Ad 3.}
Dla przedziału ograniczonego $I$ o końcach $a\leq b$ pokrycie nim samym
daje $\lambda^{\ast}(I)\leq b-a$.
Dla odcinka domkniętego $[a,b]$ powtarzamy argument
z uwagi~\ref{uw:dlugosc-przedzialu-vitali}:
dowolne pokrycie rozszerzamy do otwartego z przyrostem sumy długości
mniejszym niż $\varepsilon$, wybieramy podpokrycie skończone
i korzystamy z tego, że suma jego długości jest co najmniej $b-a$.
Z dowolności $\varepsilon>0$ wynika
$\lambda^{\ast}([a,b])\geq b-a$.

Jeżeli $a<b$, to dla $0<\delta<(b-a)/2$ odcinek
$[a+\delta,b-\delta]$ zawiera się w $I$. Monotoniczność daje
\[
  b-a-2\delta\leq\lambda^{\ast}(I)\leq b-a.
\]
Przejście z $\delta$ do zera dowodzi równości dla wszystkich rodzajów
końców. Przypadek jednopunktowy wynika z punktu 2.
Przedział nieograniczony zawiera odcinki domknięte o dowolnie dużej
długości, zatem jego miara zewnętrzna jest nieskończona.

\emph{Ad 4.}
Ustalmy $b\in\R$, $E=(-\infty,b)$ i dowolny $A\subseteq\R$.
Oznaczmy
\[
  A_1=A\cap(-\infty,b),\qquad A_2=A\cap[b,+\infty).
\]
Podaddytywność daje
$\lambda^{\ast}(A)\leq\lambda^{\ast}(A_1)+\lambda^{\ast}(A_2)$.
Wystarczy wykazać przeciwną nierówność.
Gdy $\lambda^{\ast}(A)=+\infty$, jest ona automatyczna.
Załóżmy więc, że $\lambda^{\ast}(A)<+\infty$, i ustalmy
$\varepsilon>0$. Na mocy uwagi~\ref{uw:rodzaje-pokryc} istnieje
pokrycie $A$ przedziałami otwartymi $I_n$ takie, że
\[
  \sum_{n=1}^{\infty}|I_n|<\lambda^{\ast}(A)+\varepsilon.
\]
Połóżmy
\[
  I_n^-=I_n\cap(-\infty,b),\qquad
  I_n^+=I_n\cap[b,+\infty).
\]
Są to przedziały lub zbiory puste, przy czym
$|I_n|=|I_n^-|+|I_n^+|$.
Rodziny $(I_n^-)$ i $(I_n^+)$ pokrywają odpowiednio $A_1$ i $A_2$.
W definicji miary zewnętrznej wolno użyć także tych przedziałów
i pustych składników, więc
\begin{align*}
  \lambda^{\ast}(A_1)+\lambda^{\ast}(A_2)
    &\leq\sum_{n=1}^{\infty}|I_n^-|+\sum_{n=1}^{\infty}|I_n^+|\\
    &=\sum_{n=1}^{\infty}\bigl(|I_n^-|+|I_n^+|\bigr)
      =\sum_{n=1}^{\infty}|I_n|\\
    &<\lambda^{\ast}(A)+\varepsilon.
\end{align*}
Z dowolności $\varepsilon>0$ otrzymujemy żądaną nierówność.
Zatem $E\in\mathcal{L}$.

\emph{Ad 5.}
Z twierdzenia~\ref{tw:generatory-borelowskie} i punktu 4 wynika
\[
  \mathcal{B}(\R)
  =\sigma\bigl(\{(-\infty,b):b\in\R\}\bigr)
  \subseteq\mathcal{L},
\]
ponieważ $\mathcal{L}$ jest sigma-ciałem zawierającym wszystkie
wskazane półproste.
\end{proof}
